\chapter{Actualizar y respaldar}\label{cap:actualizar}

\section{Copia de seguridad de un poste}

Haz una copia de cada poste al entregarlo y después de cada cambio importante. Guárdala con el número del poste.

\begin{enumerate}
\item Une el computador a la red del poste y abre \texttt{http://\dato{redPropiaIp}/settings}.
\item \textbf{Security \& Updates} › «Backup \& Restore»:
\begin{itemize}
\item \textbf{Backup configuration} descarga \texttt{cfg.json} (pines, LED, red, horario).
\item \textbf{Backup presets} descarga \texttt{presets.json} (escenas del poste, modo automático, escena al encender).
\end{itemize}
\item Renómbralos, por ejemplo \texttt{poste-07\_cfg.json} y \texttt{poste-07\_presets.json}.
\end{enumerate}

\nota{\texttt{cfg.json} no incluye las contraseñas: WLED las guarda aparte. Al restaurar en una placa nueva hay que volver a poner la contraseña de la red propia.}

\section{Restaurar}

En la misma pantalla, \textbf{Restore configuration} sube \texttt{cfg.json} y \textbf{Restore presets} sube \texttt{presets.json}. Después corta la corriente y vuelve a conectarla. Para reemplazar la ESP32 de un poste: graba el firmware en la nueva (capítulo~\ref{cap:grabar}), restaura los dos archivos, cambia la contraseña e instala LUMA.

\section{Actualizar LUMA}

\begin{enumerate}
\item En el computador, trae la última versión: \texttt{git pull} en \texttt{luma-app}.
\item \texttt{npm install} (por si cambiaron dependencias) y \texttt{npm test}.
\item Une el computador a la red del poste y ejecuta:
\begin{comando}
npm run build:poste
npm run instalar:poste -- 4.3.2.1
\end{comando}
\end{enumerate}
Las escenas del poste no se tocan: están en \texttt{presets.json}, no en la app.

\section{Actualizar el firmware por Wi-Fi}

\advertencia{Actualiza el firmware solo a una versión probada en un poste de banco. Haz antes la copia de seguridad.}

\begin{enumerate}
\item Si las actualizaciones están bloqueadas, desbloquéalas: \textbf{Security \& Updates}, desmarca «Lock wireless (OTA) software update», escribe la contraseña de bloqueo en «Passphrase» y guarda.
\item Abre \texttt{http://\dato{redPropiaIp}/update}.
\item Elige \texttt{WLED\_\dato{wledVersion}\_ESP32-S3\_N16R8\_actualizacion.bin} (el archivo de actualización, \textbf{no} el completo) y pulsa \textbf{Upload}.
\item Espera a que el poste se reinicie (hasta un minuto). No cortes la corriente mientras tanto.
\item Comprueba que LUMA abre y que las dos caras encienden. Vuelve a bloquear las actualizaciones.
\end{enumerate}

\nota{Si una actualización por Wi-Fi falla y el poste no arranca, graba el binario completo por USB (capítulo~\ref{cap:grabar}) y restaura la copia de seguridad.}

\section{Asegurar un poste instalado}

En \textbf{Security \& Updates}:
\begin{enumerate}
\item \textbf{Settings PIN}: un PIN de 4 cifras para los ajustes. Anótalo en un lugar seguro: si se olvida, hay que regrabar el poste.
\item \textbf{Lock wireless (OTA) software update} con una contraseña en «Passphrase». Mientras está bloqueado, nadie puede cambiar el firmware por Wi-Fi ni usar el editor de archivos \texttt{/edit}.
\item Guarda. En LUMA, escribe el PIN en «Mi montaje» › «Seguridad».
\end{enumerate}

\advertencia{La contraseña de bloqueo de actualizaciones y el PIN no se recuperan. Guárdalos con los datos del poste, por un canal privado, nunca en la hoja de verificación.}
